\documentclass{article}
\usepackage[margin=1in]{geometry}
\usepackage{longtable}
\begin{document}

\section*{Phase 8 Task 5 - Final Handoff and Closeout Package}

\textbf{Source of truth:} \texttt{Documentation/SRS.tex}, the previously generated Phase 8 Task 1 through Task 4 artifacts, and the generated closeout manifest under \texttt{reports/phase8/final\_closeout/}.\\
\textbf{Task goal:} Produce one final handoff set for demo, thesis, or defense use: a consolidated runbook path, one architecture and methodology summary, one list of key figures and tables, reproducibility references, and one final closeout memo answering whether the repo is actually v1-complete under the SRS.

\section*{Direct Answer}
\textbf{Is the project v1-complete according to the SRS?} \textbf{No, not in the current workspace snapshot.}

The repo is now much easier to defend honestly because Phase 8 has already done four important things:
\begin{itemize}
\item it classified the canonical artifact path,
\item it froze one exact reference-window reproducibility chain,
\item it formalized the canonical v1 operating mode,
\item and it wrote the final metrics and stop-condition review.
\end{itemize}

What remains missing is not documentation structure. What remains missing is \textbf{materialized runtime evidence}. The current workspace still lacks the persisted alert, evaluation, replay, backtest, and modern ML artifacts needed to satisfy the SRS completion contract fully.

\section*{Consolidated Handoff Path}
\textbf{Read order for demo / thesis / defense}
\begin{enumerate}
\item \texttt{README.md}
\item \texttt{Documentation/INDEX.tex}
\item \texttt{Documentation/SRS.tex}
\item \texttt{Documentation/phases/phase8\_canonical\_inventory.tex}
\item \texttt{Documentation/phases/phase8\_reference\_path.tex}
\item \texttt{Documentation/phases/phase8\_v1\_operating\_mode.tex}
\item \texttt{Documentation/phases/phase8\_metrics\_review.tex}
\item \texttt{Documentation/phases/phase8\_final\_closeout.tex}
\end{enumerate}

\textbf{Operator runbook path}
\begin{longtable}{p{0.28\linewidth}p{0.22\linewidth}p{0.38\linewidth}}
\textbf{Runbook} & \textbf{Classification} & \textbf{Why it stays in the handoff path} \\
\texttt{database/POSTGRES\_LOCAL\_RUNBOOK.md} & canonical & canonical bootstrap for the local PostgreSQL-backed collector and archive path \\
\texttt{database/PHASE3\_LOCAL\_RUNBOOK.md} & canonical & live detector setup and Gate 3 candidate production \\
\texttt{database/PHASE4\_LOCAL\_RUNBOOK.md} & supporting active & evidence, alerts, delivery, and analyst workflow still live here even though the runbook retains historical local-machine paths \\
\texttt{database/PHASE5\_PERSON1\_RUNBOOK.md} & supporting active & replay, validation, and backfill flow remain here until a single-owner replacement exists \\
\texttt{database/PHASE6\_PERSON1\_RUNBOOK.md} & supporting active & shadow-model registry and scoring flow remain here until a single-owner replacement exists \\
\texttt{database/PHASE6\_PERSON2\_RUNBOOK.md} & supporting active & starter-ranker training and evaluation flow remain here until a single-owner replacement exists \\
\end{longtable}

Phase 8 does \textbf{not} hide these historical filenames. Instead, it makes them explicit in the canonical handoff path so a reviewer knows exactly which runtime paths still matter.

\section*{Architecture Summary}
\textbf{Canonical v1 system story}
\begin{itemize}
\item The collector and archive layer provide the local-first replayable system of record.
\item Phase 3 converts replayable inputs into deterministic candidate episodes.
\item Phase 4 turns candidates into evidence-backed alerts with suppression and analyst workflow.
\item Phase 5 is the validation layer: replay integrity, conservative paper-trading, and holdout-style evaluation.
\item Phase 6 remains in the canonical story, but only as \textbf{shadow ML}: registry control, shadow scoring, calibration, and evaluation.
\item Phase 7 remains the advanced research layer: graph features, observability studies, ablations, thesis packaging, and scale-up ideas.
\end{itemize}

\textbf{Canonical v1 operating mode}
\begin{itemize}
\item \texttt{rule\_based\_plus\_shadow\_ml}
\item Rule-based candidate and alert logic stay authoritative.
\item ML may score, log, calibrate, and inform future promotion decisions, but it does not own the canonical alert decision path.
\item Phase 7 research outputs do not become deployment authority merely because they exist.
\end{itemize}

\section*{Methodology Summary}
\begin{itemize}
\item Preserve a replayable local-first data plane.
\item Use exact historical windows and deterministic replay whenever strong claims are made.
\item Evaluate alert quality in the SRS priority order: usefulness/precision first, then calibration, then lead time, then conservative paper-trade edge.
\item Keep ML promotion disciplined: shadow-first, reversible, baseline-compared, and calibration-aware.
\item Treat Phase 7 as thesis or research depth unless a later explicit promotion decision says otherwise.
\end{itemize}

\section*{Key Figures and Tables for Handoff}
\begin{longtable}{p{0.28\linewidth}p{0.16\linewidth}p{0.38\linewidth}}
\textbf{Artifact} & \textbf{Status} & \textbf{Use in final handoff} \\
Frozen reference chain table in \texttt{Documentation/phases/phase8\_reference\_path.tex} & available & defend the exact raw archive to research-package chain and explain which stages still lack runtime payloads \\
Operating-mode decision table in \texttt{Documentation/phases/phase8\_v1\_operating\_mode.tex} & available & explain why canonical v1 is rule-based plus shadow ML \\
Metrics bundle table in \texttt{Documentation/phases/phase8\_metrics\_review.tex} & available & show the final SRS metric order and current evidence status \\
Stop-condition ledger in \texttt{Documentation/phases/phase8\_metrics\_review.tex} & available & make the remaining risks and pause conditions explicit \\
Architecture and methodology summary in this document & available & give a reviewer one compact story across the whole repo \\
Phase 7 thesis-grade figures and ablation tables under \texttt{reports/phase7/} & not materialized in workspace & useful later for thesis depth, but not available in the current local snapshot \\
\end{longtable}

\section*{Reproducibility References}
Task 5 extends the Phase 8 generated artifact chain. The core machine-readable records are now:
\begin{itemize}
\item \texttt{reports/phase8/reference\_window\_freeze/phase8\_reference\_window\_manifest.json}
\item \texttt{reports/phase8/operating\_mode/phase8\_v1\_operating\_mode\_manifest.json}
\item \texttt{reports/phase8/metrics\_review/phase8\_metrics\_review\_manifest.json}
\item \texttt{reports/phase8/final\_closeout/phase8\_final\_closeout\_manifest.json}
\end{itemize}

These can be regenerated with:
\begin{itemize}
\item \texttt{python run\_phase8\_freeze\_reference\_path.py}
\item \texttt{python run\_phase8\_decide\_operating\_mode.py}
\item \texttt{python run\_phase8\_metrics\_review.py}
\item \texttt{python run\_phase8\_build\_closeout\_package.py}
\end{itemize}

Together, they define:
\begin{itemize}
\item the exact frozen reference window,
\item the code and documentation provenance for the end-to-end chain,
\item the canonical v1 promotion boundary,
\item the current metric and stop-condition truth state,
\item and the final Phase 8 handoff answer.
\end{itemize}

\section*{Final Closeout Memo}
\textbf{Why the answer is still ``not yet''}
\begin{enumerate}
\item The frozen reference chain is strong at the provenance level, but it is still missing later-stage runtime outputs in the current workspace.
\item The current local database snapshot has zero rows for the alert, evaluation, and backtest tables that matter for final claims.
\item The SRS requires one conservative backtest and one paper-trading evaluation; those artifacts are not materially present here.
\item The SRS requires one LightGBM or CatBoost ranker evaluated against wallet-unaware baselines; the committed Phase 6 implementation is still a linear starter ranker and the workspace does not contain the required evaluation artifact.
\end{enumerate}

\textbf{What is defensible right now}
\begin{itemize}
\item the repo has one clear canonical documentation and handoff path,
\item the project can explain its actual v1 operating mode without ambiguity,
\item the SRS metric order and stop conditions are now explicit and auditable,
\item and the remaining gaps are stated plainly instead of being hidden behind phase-local claims.
\end{itemize}

\textbf{What remains intentionally out of scope}
\begin{itemize}
\item live capital deployment or autonomous execution,
\item cloud-first migration or HA redesign before the single-instance path is fully proven,
\item automatic promotion of Phase 7 research models into canonical alert authority,
\item presenting thesis-grade research figures as if they were already required operational v1 artifacts.
\end{itemize}

\textbf{Best final reading of the repo state}
\begin{itemize}
\item The project is \textbf{defensible as a serious, structured, local-first system with a rule-based canonical v1 story and shadow ML support}.
\item The project is \textbf{not yet defensible as fully SRS-complete v1} in this workspace because the final runtime evidence packet is still missing.
\end{itemize}

\section*{Generated Artifact}
Task 5 also writes:
\begin{itemize}
\item \texttt{reports/phase8/final\_closeout/phase8\_final\_closeout\_manifest.json}
\item \texttt{reports/phase8/final\_closeout/phase8\_final\_closeout\_summary.md}
\end{itemize}

These files are the final machine-readable and human-readable handoff packet for Phase 8 closeout.

\end{document}
